\originalsection{18.315：作业6（原卷 Fall 2005）}
本组收录原题1的四小问与原题4的五小问, 共九个小问单元.
原题1(c)(d)的小阶例外及题4(d)的参数错误分别在解答中注明.

\begin{exercise}\label{mit:18315-hw6-q1}
\textbf{18.315, 原卷 Fall 2005, 作业6, 题1.}
令 $P_n$ 为全部 $n\times n$ 非负实矩阵组成的多面体,
各行、各列的元素和都为1; 令 $\Gamma_n$ 为它的顶点与边组成的图.
\begin{enumerate}[label=(\alph*)]
\item 证明 $P_n$ 的顶点对应于 $S_n$ 中的置换.
\item 证明两个置换顶点相邻当且仅当它们的相对置换只有一个非平凡循环.
\item 由(b)推出 $\Gamma_n$ 的直径为2, 并检查小阶情形.
\item 证明 $\Gamma_n$ 含有 Hamilton 圈, 并检查小阶情形.
\end{enumerate}
\end{exercise}
\sourceof{官方 HW6 第1页题1(a)--(d); 编者独立解答, 非官方答案;
“一个循环”排除恒等置换, (c)须 $n\ge4$, (d)须 $n\ge3$}
\begin{solution}
(a) 对 $X=(x_{ij})\in P_n$, 建立行点 $r_i$ 与列点 $c_j$ 的二部支持图,
$r_ic_j$ 是边当且仅当 $x_{ij}>0$.
对任意行点集 $I$, 非零元素只能落在邻列 $N(I)$ 上, 故
\[
|I|=\sum_{i\in I}\sum_jx_{ij}
\le\sum_{j\in N(I)}\sum_ix_{ij}=|N(I)|.
\]
Hall 定理给出一个完美匹配, 对应置换矩阵 $M$.
令 $t$ 为这些匹配位置上的最小元素. 若 $t=1$, 行和条件迫使 $X=M$;
若 $0<t<1$, 则 $Y=(X-tM)/(1-t)$ 仍非负, 行列和仍为1, 且比 $X$ 少至少一个正位置.
因此 $X=tM+(1-t)Y$. 对正位置数归纳, 得到每个 $X$ 都是置换矩阵的凸组合.

每个置换矩阵又确为顶点. 若 $M=tX+(1-t)Y$, $0<t<1$,
则 $M$ 的每个零位置由非负性迫使 $X,Y$ 同样为零;
每行只剩 $M$ 的那个非零位置, 行和条件迫使该元素为1, 所以 $X=Y=M$.
另一方面, 若一个点是多个不同置换矩阵的非平凡凸组合, 就不是顶点.
这完整证明了顶点恰为置换矩阵.

(b) 同时重排列不会改变多面体或它的边, 因而只须比较恒等矩阵与 $M_\sigma$.
把两套匹配取并. 相同匹配边独立成为一个二点分支;
其余分支是长度至少4的交替偶圈, 恰对应 $\sigma$ 的非平凡循环.
假设这类圈有 $k$ 个. 限制矩阵在并集之外的坐标全为零,
得到 $P_n$ 的一个面: 每个坐标均非负, 令坐标为零就是取一个支持超平面,
有限次相交仍为面.

在一个交替圈上, 行和、列和条件逐个迫使元素沿圈交替为 $t_j,1-t_j$,
$0\le t_j\le1$; 不同圈之间没有耦合, 相同匹配边上的元素固定为1.
因此这个面就是 $k$ 维立方体 $[0,1]^k$, 两个原顶点是相反的角.
若 $k=1$, 该面是一条边, 两顶点相邻.
若 $k\ge2$, 两点的中点还可表示为另外两个角的平均:
在第一圈选第一套匹配、其余圈选第二套得到 $R$,
反过来得到 $S$, 则 $(I+M_\sigma)/2=(R+S)/2$.
$R,S$ 不在原对角线段上. 若该线段是面, 面的定义将迫使 $R,S$ 也在线段上,
矛盾. 所以不是边. 这给出要求的充要条件.

(c) 以下置换从右向左复合. 对任意非恒等 $\sigma$, 写出其 $r$ 个非平凡循环,
每个循环选一点 $a_i$. 若 $r=1$, 它与恒等置换已经相邻.
若 $r\ge2$, 令 $\tau=(a_1\ a_2\ \cdots\ a_r)$.
$\sigma\tau$ 把各循环接成一个循环: 从 $a_i$ 先跳到下一循环的
$\sigma(a_{i+1})$, 沿该循环走回 $a_{i+1}$, 再跳到下一个循环;
依次经过全部 $r$ 个循环的点而返回起点, 其余点固定.
所以恒等置换与 $\sigma\tau$ 相邻;
而 $(\sigma\tau)^{-1}\sigma=\tau^{-1}$ 也只有一个非平凡循环,
故 $\sigma\tau$ 与 $\sigma$ 相邻. 左乘任意置换是图的自同构,
因此任意两顶点距离至多2.

当 $n\ge4$ 时, 恒等置换与 $(1\ 2)(3\ 4)$ 不相邻, 距离恰为2,
于是直径恰为2. 当 $n=1$ 时图只有一点, 直径0;
$n=2$ 时图为 $K_2$, 直径1;
$n=3$ 时每个非恒等置换都是一个换位或三循环, 图为 $K_6$, 直径1.
原题的“直径为2”因此必须加上 $n\ge4$.

(d) 构造全部置换的循环次序, 使相邻两项仅交换相邻位置的两个字母.
这种交换是换位, 由(b)必然给出 $\Gamma_n$ 的边.
$n=3$ 时可取
\[
123,\quad213,\quad231,\quad321,\quad312,\quad132,\quad123.
\]
假定 $n-1$ 个字母已经排成这样的 Hamilton 循环.
在其中每一个置换的块里, 把新字母 $n$ 插入所有 $n$ 个位置:
第一块从最右位置一路移到最左位置, 第二块从最左一路移到最右,
此后交替. 块内每一步交换 $n$ 与相邻字母;
相邻块的共同边界中 $n$ 都在同一端, 删去 $n$ 后恰是归纳循环中相邻的两项,
其一次相邻交换在插入端点字母后仍为一次相邻交换.
共有 $(n-1)!$ 个块; 当 $n\ge4$ 时它是偶数,
最后一块的末项与第一块的首项都把 $n$ 放在最右端,
因此最后与最初同样相邻.
每个 $n$ 字母置换有唯一的删去 $n$ 后的置换与插入位置,
故各块没有重复, 且恰覆盖全部 $n!$ 个置换. 归纳完成.
$n=1,2$ 时 $\Gamma_n$ 分别为 $K_1,K_2$,
按本书长度至少3的简单圈约定都没有 Hamilton 圈, 是原题的小阶例外.
\end{solution}

\begin{exercise}\label{mit:18315-hw6-q4}
\textbf{18.315, 原卷 Fall 2005, 作业6, 题4.}
把图 $G$ 的全部顶点放在一条直线 $L$ 上, 边放在若干以 $L$ 为公共边界的半平面中;
同一半平面中的边内部互不相交. 最少的半平面数称为书厚度 $c(G)$.
\begin{enumerate}[label=(\alph*)]
\item 可平面 Hamilton 图满足 $c(G)\le2$.
\item 不可平面图满足 $c(G)\ge3$.
\item $c(K_n)=n/2+O(1)$.
\item 原卷印作 $c_n(H_n)=\theta(\log n)$, $H_n$ 为 $n$ 维超立方体;
判定其正确性, 并给出可证明的修订.
\item 可平面图满足 $c(G)\le1000$.
\end{enumerate}
\end{exercise}
\sourceof{官方 HW6 第1页题4(a)--(e); 编者独立解答, 非官方答案;
原题(d)的 $c_n$ 未定义, 按书厚度 $c$ 理解后结论仍有参数笔误}
\begin{solution}
先固定一个顶点线性次序. 两条无公共端点的边能够在同一页上无交叉绘制,
当且仅当其端点不交替, 即不存在次序 $a<b<c<d$ 及边 $ac,bd$.
必要性来自 Jordan 分离: 边 $ac$ 与脊线上从 $a$ 到 $c$ 的线段围成一条闭曲线,
另一条边的两个端点分处它在半平面的两侧.
充分性则把每条边画为以两端点为直径端点的半圆;
不交替的两个区间或者内部分离、或者一个包含另一个,
对应半圆内部便不会相交. 共端点的半圆也只有端点相交.
以下均可用这一判据来检查页面.

(a) 在一张平面嵌入中取 Hamilton 圈. 按圈上的循环顺序把顶点排在线上.
其余边处于圈的内侧或外侧; 同侧的两条边端点不可能交替,
否则 Jordan 分离迫使它们相交. 因此两侧各分配一页即可.
圈边连接线性次序中的连续顶点, 或最初与最后顶点,
不会与任何边端点交替, 可全部放在第一页. 故 $c(G)\le2$.

(b) 若至多两页, 把第一页放在脊线上方、第二页放在下方,
就得到一张平面嵌入; 页内互不相交, 两页内部又位于不相交的半平面.
所以不可平面图不能有这样的两页画法, 必须 $c(G)\ge3$.

(c) 实际上 $n\ge4$ 时有 $c(K_n)=\lceil n/2\rceil$.
为证上界, 将顶点依次编号 $0,1,\ldots,n-1$,
把边 $ij$ 依 $i+j\pmod n$ 分组, 每页放两个连续的余数类:
$\{0,1\},\{2,3\},\ldots$, 若余数个数为奇数则最后一类单独成页.
若 $a<b<c<d$ 的交替边 $ac,bd$ 在同一页, 则
\[
2\le(b+d)-(a+c)=(b-a)+(d-c)\le n-2.
\]
它们的和模 $n$ 因而不能相同或相差1, 与分组矛盾.
每页没有交替边, 页数为 $\lceil n/2\rceil$.

为证下界, 暂把线性次序首尾相接成圆周.
每一页所含的对角线仍不交叉; 再补上全部 $n$ 条圆周边,
得到一个所有顶点都在外面边界上的简单平面图.
加对角线直到内面全为三角形, 每次加边保持无交叉.
若最终有 $d$ 条对角线, 则边数 $n+d$、内面数 $d+1$;
数内面边界, 每条对角线两次、圆周边一次, 得
$3(d+1)=2d+n$, 即 $d=n-3$.
因此一页至多含 $n-3$ 条 $K_n$ 的圆周对角线,
而整图有 $\binom n2-n=n(n-3)/2$ 条, 故页数至少 $n/2$.
取整数上界即得到等号. 小阶 $K_1$ 无边可取0页,
$K_2,K_3$ 各需1页, 都与所要求的渐近式一致.

(d) 写 $Q_n$ 为 $n$ 维超立方体, 顶点是二进制串, 相差恰一位时相邻.
它有 $N=2^n$ 个顶点、$M=n2^{n-1}$ 条边.
任何一页都是外平面图; 由上一问补圆周边、三角剖分的计数,
当 $N\ge3$ 时一页至多有 $2N-3$ 条边. 因而 $n\ge2$ 时
\[
c(Q_n)\ge\frac{n2^{n-1}}{2^{n+1}-3}>\frac n4.
\]
这已经否证原卷的 $\Theta(\log n)$: $n/\log n$ 无界.

上界也可自足构造. $Q_1$ 一条边需一页.
从 $Q_{n-1}$ 的 $n-1$ 页画法出发, 将两个副本依次放在脊线上,
第一个沿原次序, 第二个沿完全相反的次序.
各副本内部的边复用原来的 $n-1$ 页; 反转次序保持不交替,
而两个副本的区间不重叠, 跨副本的内部边也互不交替.
最后把对应顶点间的匹配放在一张新页: 若第一副本次序为
$v_1,\ldots,v_{2^{n-1}}$, 第二副本为
$v'_{2^{n-1}},\ldots,v'_1$,
匹配边 $v_iv'_i$ 的区间两两嵌套, 所以不交叉.
归纳得 $c(Q_n)\le n$. 合并上下界,
\[
\frac n4<c(Q_n)\le n\quad(n\ge2),\qquad c(Q_n)=\Theta(n).
\]
若改用顶点数 $N=2^n$ 作为参数, 才能写成 $\Theta(\log N)$.
这是一项明确修订, 不是对原题文字的无声替换.

(e) 我们自足证明更强的统一界 $c(G)\le5$, 所以必有 $c(G)\le1000$.
以上已经证明同页没有端点交替就可以画成不交半圆, 下面只需构造
同一个全局顶点次序和五种页面颜色.

先把给定简单平面图扩充成一张简单平面三角剖分.
少于三个顶点的图直接用至多一页; 以下至少三个顶点.
在原嵌入中不断加可以无交叉画出的非重复边, 直到不能再加.
这种过程有限, 因为顶点集固定而候选边有限.
所得图连通, 否则同一个面的两个分支间还可加边.
它没有割点: 若 $v$ 是割点, 在其转角次序中取邻接不同删除分支的
两条相继边 $vu,vw$, 沿两边附近绕开 $v$ 加入 $uw$.
$u,w$ 之间本无边, 否则它们在删去 $v$ 后已经相连, 而这条新边没有交叉.
连通且没有割点的至少三点平面图, 每个面边界都是简单圈:
若某面边界重复经过 $v$, 可在该面中从 $v$ 画出一条小分隔弧,
圈的两次到达之间的边分处它的两侧, 删去 $v$ 就使这两部分分离,
与没有割点矛盾.
若一个面边界长度至少4, 取其四个相继顶点 $v_1,v_2,v_3,v_4$.
边 $v_1v_3,v_2v_4$ 若都已存在, 只能落在此面补集的另一侧盘中,
但其端点交替, Jordan 分离迫使两边相交. 因此至少一条不存在,
可在面内加入, 与极大性矛盾. 所以所有面都是三角形.
下面把这张固定初始三角剖分的顶点称为\textbf{基础顶点};
构造中再增加的辅助顶点及边在最后删除, 所得画法限制回原图.

\medskip
\noindent\textbf{一、两层盘的平面次序。}
设 $K=v_1v_2\cdots v_kv_1$ 是一个无弦的外圈，圈内已有一个连通的
外平面图 $I$。每个 $I$ 的顶点都在 $I$ 的外面边界上，并邻接至少一个
$v_i$；$K$ 与 $I$ 之间的所有面都是三角形。
$I$ 的内部可以有更深的顶点，但在本段将它们及其边删去，只保留
外圈边、$I$ 的边及连接 $K$ 与 $I$ 的\textbf{连接边}。
任取顺时针方向，把 $K$ 按此方向编号。

包含外圈边 $v_kv_1$ 的内侧三角形的第三个顶点记为 $a$。
在这个三角形内增加一个顶点 $a_0$，接到 $v_1,v_k,a$。
这样 $a_0$ 在 $I$ 中只邻接 $a$，以 $a_0a$ 为根块。
此辅助操作始终可做；以下把操作后的图仍记为 $I$。

$I$ 的\textbf{块}是它的极大二连通子图，桥也单独算一个块。
一个非平凡块的外面边界是一个简单圈，其余边都是圈内的弦；
平凡块则只有一条边。以 $a_0$ 为根，从 $a_0$ 到某个块 $B$ 的任一
内层路径第一次进入 $B$ 的顶点相同，记为 $u(B)$，称为该块的
\textbf{首顶点}。根块的首顶点为 $a_0$。
非根块进入前的最后一条边所属的块是它的父块。这样得到一个有根块树：
一个割点有多个后续块时，这些块有同一个父块。把一个内层顶点分配给
包含它的最靠近根的块。根块分配到自己的全部顶点；每个非根块分配到
自己的全部顶点，首顶点除外。

这些块分解事实也可以从路径直接看出。若一个块能从根经两个不同顶点
进入，则两条进入路径与块内路径组成一个圈，把块和沿途的边放入同一个
二连通子图，违背块的极大性。反复删去末端块，得到树结构以及上述唯一
的顶点分配。

对一个内层顶点 $x$，令 $\Gamma(x)$ 为它在外圈上的邻点集合。
对块 $B$，令 $\Gamma(B)$ 为\textbf{分配给 $B$ 的顶点}的外圈邻点集合。
集合 $S\preceq T$ 的意思是 $S$ 的每个编号都不大于 $T$ 的每个编号。
设一个块的顺时针边界次序为
\[
 u,u_1,\ldots,u_t,u,\qquad u=u(B);
\]
桥的情形取 $t=1$。有如下平面分离事实。
\begin{enumerate}
\item $\Gamma(u_1)\preceq\Gamma(u_2)\preceq\cdots\preceq\Gamma(u_t)$。
\item 若 $f$ 是 $\Gamma(u_1)$ 的最小编号，$l$ 是
  $\Gamma(u_t)$ 的最大编号，则 $f<l$。
  $v_fuu_1$ 与 $v_lu_tu$ 是外侧三角面；$u$ 不邻接编号严格位于
  $f,l$ 之间的外圈顶点。
\item 若 $B'$ 是 $B$ 的祖先，则
  $\min\Gamma(B')\leq \min\Gamma(B)$ 且
  $\max\Gamma(B')\geq \max\Gamma(B)$。
  若 $B'$ 不是 $B$ 的后代，也不是 $B$ 本身，则 $B'$ 不邻接
  编号严格位于 $f,l$ 之间的外圈顶点。
\item 若 $B,B'$ 互不为祖先，则
  $\Gamma(B)\preceq\Gamma(B')$ 或
  $\Gamma(B')\preceq\Gamma(B)$。
\end{enumerate}

为完整说明其中的平面论证，取从 $a_0$ 到 $u$ 的一条简单路径 $R$，
其内部与 $B$ 不交；这是块树给出的根路径。在外圈内，沿
$v_1a_0$、$R$ 及 $a_0v_k$ 切开一条任意窄的带，并删去 $B$ 的内部
（桥则沿桥的两侧切开）。在块外、外圈内的区域因而成为一个盘。
它的两条相对边界上，外圈顶点依次为 $v_1,\ldots,v_k$，
块的非首顶点依次为 $u_t,\ldots,u_1$。
原来从 $u_i$ 到外圈的连接边成为盘内的弧，不穿越切带。
盘内两条弧的端点若沿盘边界交替出现，两弧必相交：
第一条弧与盘的一段边界构成 Jordan 曲线，第二条弧的两个端点在其两侧。
因此 $i<j$ 时，$u_i$ 的外邻点不可能在 $u_j$ 的外邻点之后，证明第1点。

边 $uu_1$ 的块外一侧属于一个三角面；第三点必在外圈上，
因为该侧在 $I$ 的外面，而所有这一侧的面都已三角化。
记第三点为 $v_s$。在切开的盘中，弧 $uv_s$ 与外圈的前段边界
隔开 $u_1$ 和所有 $v_i$（$i<s$）；若 $u_1$ 邻接这些点，连接边
必须跨越这条弧。故 $s=\min\Gamma(u_1)=f$。
对末边 $u_tu$ 作镜像论证，第三点正是
$v_l$，$l=\max\Gamma(u_t)$。
两面所用的外顶点不同。桥的两侧若用同一个第三点，就会把同一三角形
的两侧都指定为面，使外圈其余顶点无处位于圈内。
若块非平凡而两面第三点都为 $v_f$, 在 $u$ 处三角面
$v_fuu_1$ 与 $v_fu_tu$ 分别占据边 $uv_f$ 的两侧;
边 $uu_1,uu_t$ 之间朝块圈外的整个转角扇区因此恰被这两个面占满,
不容任何额外内层边进入。根路径 $R$ 的最后一条边不属于 $B$,
也不可能是接到外点 $v_f$ 的连接边, 故只能从朝块圈内部的扇区进入 $u$。
但 $a_0$ 在块圈外, $R$ 的内部与 $B$ 不交; Jordan 分离使 $R$ 无法
进入块圈内部, 矛盾。于是 $f<l$。
两条弧 $uv_f,uv_l$ 与外圈从 $v_f$
到 $v_l$ 的前向弧围成一个区域，$u_1,\ldots,u_t$ 朝该区域一侧，
根路径在另一侧。在 $u$ 处, 边 $uv_f,uu_1$ 因三角面而连续,
$uu_t,uv_l$ 也连续。若块非平凡, 朝该区域的中间转角扇区是块圈内部,
一条从 $u$ 到开外弧 $(v_f,v_l)$ 的边必须从此扇区出发而穿过块圈,
或从另一扇区出发而跨越边 $uv_f,uv_l$; 两者都不可能。
若块为桥, 这两个面在桥的两侧, 朝开外弧的转角扇区被它们占据,
同样没有新的连接边可通过。故 $u$ 没有中间外邻点, 证明第2点。

父块含有 $u$，而且 $u$ 分配给父块。第2点给出它同时邻接
$v_f,v_l$；逐层向根推进，得到第3点的两个端点不等式。
还可用刚才的切盘把外圈内分为三个区域：块的前侧区域、后侧区域，
以及朝外圈开弧 $(v_f,v_l)$ 的中间区域。任何内层路径从根进入中间
区域，都必须先经过 $B$；否则会穿过上述 Jordan 边界。
因此能邻接开弧 $(v_f,v_l)$ 的其他块只能是 $B$ 的后代。
不与 $B$ 有祖先关系的块，其分配顶点诱导一个连通子图，且不能穿过
$B$ 或根路径；它只能全在前侧或全在后侧。这两侧分别只能到达
不大于 $f$、不小于 $l$ 的外顶点，证明第3点的其余断言及第4点。
这些论证同样适用于桥：把桥的两侧看作上述两条块边界即可。

令
\[
 d(B)=\min\Gamma(B),\qquad r(B)=\max\Gamma(B),
\]
称 $v_{d(B)}$ 为块的\textbf{支配顶点}。
对非根块有 $d(B)=f,r(B)=l$；根块是 $a_0a$，
而 $a$ 邻接 $v_1,v_k$，故根块也有 $d=1,r=k$。
每块均有 $d(B)<r(B)$，特别是没有块被 $v_k$ 支配。
祖先关系中 $d$ 不减。两个互不为祖先的块不能有相同的 $d$，
否则第4点及 $d<r$ 相矛盾。于是同一个外顶点所支配的块构成块树中
一条向下的连续路径：若两个同支配块之间有第三个块，单调性迫使
第三块的 $d$ 也相同。

还需要明确一个涉及割点的细节。设 $q$ 是使
$v_{d(B)}$ 邻接 $u_q$ 的最大下标。由第1点，
\[
 \Gamma(u_i)=\{v_{d(B)}\}\qquad(1\leq i<q).
\]
这些 $u_i$ 不属于其他块。否则它们必须是某个后续块的首顶点，
但第2点要求任何首顶点至少邻接两个不同外顶点。
如果 $B$ 有一个同支配的子块，则这个子块的首顶点只能是 $u_q$：
它不能是 $u_i$（$i<q$），不能是 $u_i$（$i>q$，不邻接该支配点），
也不能是 $u$，因为通过 $u$ 到达的其他块与 $B$ 同父，而非 $B$ 的子块。
根首顶点 $a_0$ 在内层的度数为1，也不可能是子块首顶点。

\medskip
\noindent\textbf{二、两层图的连续块三页布局。}
三页记为 $P,A,B$。外圈依次排成 $v_1,\ldots,v_k$。
在 $v_i,v_{i+1}$ 之间，依块树从上到下的顺序放入全部由 $v_i$
支配的块。对每个块，分配顶点沿块边界\textbf{逆时针}排列；
非根块的次序为 $u_t,u_{t-1},\ldots,u_1$，根块的次序为 $a_0,a$。
同一块分配到的顶点是一个连续段。块的首顶点总在其他顶点之前：
它分配给父块，而父块的支配编号较小，或支配相同而先行。

把块树以 $A,B$ 二染色，相邻的父子块颜色相反。
每条内层边取它所属块的颜色。
连接边 $v_ix$ 若 $v_i$ 位于 $x$ 之前，称为\textbf{后向连接边}，
并给颜色 $P$；否则称为\textbf{前向连接边}，给与 $x$ 所分配块
相反的 $A$ 或 $B$。外圈边给颜色 $P$。
这里“前向”“后向”只描述上述端点次序，并不赋予图有向结构。

首先，后向连接边只能从块的支配点出发。它的另一个端点在该外点
紧接的插入段里，故所有后向边与外圈边都不交错。
同一块的内层边也不交错，因为该块的顶点的相对次序正是外面边界的
循环次序，交错弦会在平面内相交。

其次，任意两条前向连接边不交错。设它们为 $xv_i,yv_j$，其中
$x<y$。要排除交错，只须证明 $v_i\leq y$ 或 $i\geq j$。
若 $x,y$ 分配给同一块且均非首顶点，逆时针排列与第1点直接给出
$i\geq j$。同块分配的首顶点只可能是根首点 $a_0$；它唯一的前向
外邻点是 $v_k$，同样有 $i\geq j$。
若两端分配给不同块 $X,Y$，有 $d(X)\leq d(Y)$；等号时 $X$ 在
同支配路径上先于 $Y$。若两块无祖先关系，第4点和 $x<y$ 给出
$\Gamma(X)\preceq\Gamma(Y)$，于是 $i\leq d(Y)$，$v_i<y$。
若有祖先关系，则 $X$ 是 $Y$ 的祖先。第3点说 $X$ 的邻点不在
$(d(Y),r(Y))$ 内。因此 $i\leq d(Y)$ 时仍有 $v_i<y$；
否则 $i\geq r(Y)\geq j$。所有情形均不交错。

最后，设一对内层边或前向连接边交错，端点为
\[
 w<y<x<z,\qquad wx, yz.
\]
边 $wx$ 不能是前向连接边。否则 $x$ 在外圈上，$yz$ 必为内层边
（两条前向边已经排除）。设 $yz$ 所在块为 $D$。外点 $x$ 位于
两个块顶点 $y,z$ 之间，而该块的非首顶点连续，故 $y$ 必为首顶点。
由第2点，$y$ 邻接 $v_{r(D)}$，且 $v_{r(D)}>z$。
于是 $wx$ 与 $yv_{r(D)}$ 是交错的两条前向边，矛盾。

故 $wx$ 是某块 $D$ 的内层边，$x$ 是其非首顶点。
还须有 $v_{d(D)}<y$。若 $y<v_{d(D)}<x$，则 $w$ 是 $D$ 的
首顶点，而且 $w$ 邻接 $v_{d(D)}$。边 $wv_{d(D)}$ 是前向连接边，
与 $yz$ 交错，违背上一段。$y$ 不可能就是支配外点，因为 $yz$
既为内层边又或为前向连接边，其左端点必须是内点。

所有位于 $v_{d(D)},x$ 之间的内点都属于同支配路径上的块。
若 $w$ 是 $D$ 的非首顶点，连续性立即给出 $y$ 分配给 $D$。
若 $w$ 是首顶点，$y$ 除分配给 $D$ 外，只可能落在父块的连续段中
且位于 $w$ 之后；更早的祖先段在 $w$ 之前，而父块与 $D$ 之间没有
其他块。在后一情形，父块和 $D$ 有同一支配点，$w$ 正是前面定义的
$u_q$；父块连续段在 $w$ 之后的顶点均为 $u_{q-1},\ldots,u_1$。
它们没有前向连接边，且全部内层边属于父块。父块的全部顶点又在
$x$ 之前，所以它们不可能有一条这里要求的边 $yz$，其中 $z>x$。
因此后一情形不发生，$y$ 必分配给 $D$。

若 $yz$ 是前向边，其颜色与 $D$ 相反；若 $yz$ 是内层边，它不属于
$D$（同块边不会交错），而 $y$ 分配给 $D$，故 $y$ 是另一个子块
$D'$ 的首顶点，$D'$ 的颜色与 $D$ 相反。这样任何交错的两边都异色，
两层三页布局得到完整证明。

此构造还给出递归所需的附加信息。对每个非平凡块 $D$，设其颜色为
$X\in\{A,B\}$，另一个颜色为 $\overline X$。
块的圈边和弦全部为 $X$；其非首顶点构成一个连续段。
接到这些非首顶点的\textbf{块外}边，只可能是颜色 $P$ 的后向边、
颜色 $\overline X$ 的前向边，或颜色 $\overline X$ 的子块边。
颜色 $X$ 的块外边若接到 $D$，只能接到首顶点。
此外，所有新边中，接到外圈首点 $v_1$ 的边都是后向边，颜色全为 $P$。

\medskip
\noindent\textbf{三、统一次序的五页递归。}
现在固定五种颜色。递归对象是一张尚未展开的三角剖分盘，其简单外圈
记为 $C=c_1c_2\cdots c_pc_1$。圈顶点已放入\textbf{同一个全局次序}，
内部的基础顶点均未放入。维持以下不变量。
\begin{enumerate}
\item 外圈在全局次序中的相对次序为 $c_1,\ldots,c_p$；
  $c_2,\ldots,c_p$ 是一段连续的旧顶点。
  $c_1,c_2$ 之间允许有其他已放置顶点，记这一段为 $N$。
\item 圈边及所有圈内的外圈弦已经着同一种颜色 $P$；特别是长边
  $c_1c_p$ 已着 $P$。整个已着色部分没有同色交错。
\item 接到 $c_2,\ldots,c_p$ 的已经着色的盘外边，至多使用另两种
  颜色 $Q,R$，且 $P,Q,R$ 两两不同。因此盘外 $P$ 边若接到外圈，
  只能接到 $c_1$。
\item 所有尚未展开盘的非首顶点段两两不交；盘可以在割点处共用首点，
  但每一个非首顶点只属于一个这样的连续段。
\end{enumerate}
外三角形的三个顶点按圈序排列，三条边都着 $P$。
它没有盘外边，任选两个异于 $P$ 的颜色作 $Q,R$；不变量成立。

若当前盘内部没有基础顶点，则不再展开：其所有原有边只有圈边或圈弦，
均已着色。否则先进行以下\textbf{仅用于本次平面构造}的预处理。
把外圈弦从当前工作的平面盘中擦去，但保留它们已经着 $P$ 的书边。
擦去后图仍二连通。原盘的每个面边界都是简单圈，故原盘二连通；
删去任意顶点 $z$ 后，$C-z$ 仍连通。所擦去的边的两端均在 $C-z$
的同一连通分支中，故它们不可能把其他分支重新连通。
既然原图减去 $z$ 连通，擦去弦后的图减去 $z$ 也连通。
因此新面的边界仍是简单圈。

每个新出现的非三角面由原三角面沿被擦去的外圈弦合并而成。
其中每个参与合并的三角面都含有一个被擦去的弦，故含有至少两个
外圈顶点；它的第三点若不在外圈，便是一个邻接外圈的基础顶点。
所以新面的边界只含外圈顶点和原先邻接外圈的基础顶点，而且含有
外圈顶点。在每个非三角面中加一个面中心，接到所有边界顶点，
把该面三角化。各中心彼此不邻接；每个中心都邻接外圈。
已擦去的旧弦不再参与这个平面图，因而不要求它们与辅助扇边同时
构成平面画法；它们在书中的颜色和端点次序始终保留。

取工作盘删除外圈后的外面边界顶点，诱导图记为 $I$。
它连通且外平面，并且这些顶点恰好是工作盘中邻接外圈的内点。
以下从邻扇逐步核查。对每个 $c_i$, 在盘内按转角顺序列出它的邻点,
从 $c_{i-1}$ 到 $c_{i+1}$; 中间点都严格位于圈内, 因其他圈点作邻点会产生外圈弦。
连续邻点与 $c_i$ 围成三角面, 所以其中的内点形成连通邻点链。
每条外圈边 $c_ic_{i+1}$ 的内侧三角面第三点 $z_i$ 必为内点:
否则该三角形有一个圈弦, 或整个盘只有无内点的三角形, 而当前盘有内点。
于是 $c_i,c_{i+1}$ 的两条邻点链共享 $z_i$, 全部外圈邻点的并连通。
工作盘的任意内点分支必含一个外圈邻点, 沿一条到外圈的路径取第一条到圈的边即可。
因此删除外圈后的整个余图连通, 这些邻点的诱导图 $I$ 也连通。
删除外圈所打开的外面恰由原来含圈点的三角面相连组成:
一个外圈邻点因其到圈的边被删除而暴露于该外面;
反过来, 暴露点位于其中一个三角面的边界, 因而邻接这个面的圈点。
所以 $I$ 的每个顶点都在其外面边界上, $I$ 是外平面图。
余下各点全在 $I$ 的某个有界面内。
每个有界面由一个非平凡块的简单圈围住，因而这些余点分别归入各
非平凡块圈的内部。圈外紧贴 $I$ 边界的三角面第三点在外圈，
所以恰好满足第一部分的两层盘条件。

加根点 $a_0$ 后，选两种不属于 $\{P,Q,R\}$ 的颜色 $A,B$，
应用上述连续块构造。新顶点的局部次序嵌入全局次序时，
对支配点 $c_i$（$i\geq2$），把其插入段放在原本空的
$(c_i,c_{i+1})$ 中；对 $c_1$，把其插入段放在 $N$ 的\textbf{后面}、
$c_2$ 的前面。所以旧顶点的次序完全不变，$N$ 中的全部旧顶点仍在
这一新插入段的左边。没有块被 $c_p$ 支配，无须在 $c_p$ 之后插入。
给内层边和连接边着色时，完全使用第二部分的规则。

必须检查新边与\textbf{全部旧边}之间没有同色交错，而不能只检查当前
平面盘。所有新 $A,B$ 边的端点都在 $N$ 的右边和 $c_p$ 之间，
并且不接到 $c_1$。如果某条旧边与这样的新边交错，则旧边至少有一个
端点严格位于新边的端点区间内。这里唯一可能的旧顶点是
$c_2,\ldots,c_p$；接到它们的旧边使用 $P,Q,R$，不用 $A,B$。
故没有旧 $A,B$ 边与新 $A,B$ 边交错。

新 $P$ 边都是后向连接边。若左端点是 $c_i$，$i\geq2$，
其右端点在原来空的 $(c_i,c_{i+1})$ 中，旧边不可能有端点严格落在
这个区间，因此不交错。
剩下的是新边 $c_1x$，其中 $N<x<c_2$，意思是 $N$ 的全部顶点
都在 $x$ 之前。设一条旧 $P$ 边的端点为 $s<t$，可能的两种交错
必须逐一排除：
\begin{itemize}
\item $s<c_1<t<x$。此时 $t\in N$，旧边 $st$ 与已经着 $P$ 的
  长边 $c_1c_p$ 交错，违背旧布局合法性。
\item $c_1<s<x<t$。此时 $s\in N$。由于 $N$ 的全部旧点均在
  $x$ 之前，$t$ 要么是某个 $c_j$（$j\geq2$），要么在 $c_p$ 之后。
  第一种是接到非首圈点的盘外 $P$ 边，违背第3不变量；
  第二种与长边 $c_1c_p$ 交错，也违背旧布局合法性。
\end{itemize}
故新旧同页均不交错。新边之间的合法性已由两层引理证明。

接下来对每个含有未放置基础顶点的非平凡块圈分别递归。
若该块颜色为 $X$，其圈边和弦都已经是 $X$；
圈的次序从首顶点起，正是局部逆时针边界次序。
它的非首顶点是连续段，其盘外边在这些顶点上只使用
$P,\overline X$，两者均异于 $X$。因此下一步的不变量可取
\[
 P_{\rm next}=X,\qquad Q_{\rm next}=P,
 \qquad R_{\rm next}=\overline X.
\]
各新连续段两两不交，也不进入此前尚未展开的其他非首点段。

还要确认先展开一个盘不会破坏另一个待展开盘的不变量。
一次展开只在自己的非首点段内部，以及该段第一点之前的紧邻空隙
插入新点；它不进入另一个待展开盘的非首点段。
不同待展开盘的内部互不相交，原图的新边若接到另一个盘的非首顶点，
该顶点只能是本盘的首顶点。共用这类边界割点的两个待展开盘必来自
同一次两层分解：更深子盘的边界全在其祖先盘的严格内部，不会再含
这个旧外圈点。这种接到共享点的新边全部取本盘圈色。
在块树中，本盘是以该点为首顶点的子块，圈色正是另一个盘外部允许的
相反块色。因此另一个盘的连续性和 $Q,R$ 限制均被保留。
更深展开只接到自己的边界，不再接这个旧外圈点。
若接到另一个盘的首顶点，则第3不变量本来没有限制该点的盘外颜色。
辅助边只接本次工作的边界和新点，同样符合上述检查。
这样可以依任意固定次序逐个展开各盘，而全程只维护一个全局顶点次序。

最后证明过程有限，排除不断增加辅助点导致无限递归的可能。
对一个待展开盘，以其严格内部的基础顶点数 $m$ 为归纳参数。
本次新增的每个面中心以及根点 $a_0$ 都邻接外圈；连接到外圈的边
不能穿越一个内层块圈，故这些新点都不可能位于内层块圈的严格内部。
之前新增的辅助点同理，已在某个祖先的第一层中，不能在本盘内部。
面中心之间无边，而 $a_0$ 在内层只属于一个桥，所以每个非平凡内层
块圈必含至少一个本轮原有的内层顶点，即一个基础顶点。
这个基础顶点被放到了子盘边界，已不在子盘严格内部。
因此每个真正递归的子盘的内部基础顶点数至多为 $m-1$。
当 $m=0$ 时停止，故归纳终止。每步分支数有限，整棵递归树有限。

所有原图边最终都着色：第一层的圈边、块弦和连接边在当步处理；
剩余边及其未放置端点均在某个非平凡块圈内，交给该子盘。
初始及后来为证明引入的全部辅助顶点最后删去。
所得原图在统一顶点次序下使用至多五页，故其书厚度至多 $5<1000$。

\medskip
\noindent\textbf{来源边界。}
连续块的两层布局参照 Yannakakis 的原论文第二节；以上把所需平面分离、
割点尾段及交错检查全部展开，并单独给出五页递归不变量及其旧边审查。
原论文中四页定理及其专门递归算法不作为本证明的前提。
研究来源是 M. Yannakakis, Embedding planar graphs in four pages,
Journal of Computer and System Sciences 38(1), 1989, 36--67,
\href{https://page.math.tu-berlin.de/~felsner/Lehre/SemInter/Yannakakis-4pages-1989.pdf}{原文链接}.
受限原论文只留链接, 不附扫描件或原图; 上述构造使用重新表述的完整证明。
\end{solution}
